\begin{center}
    \vspace*{1.5em}
    {\LARGE\heiti\bfseries 跨模态抗污染异常检测算法对比与主动防御基准评估 \par}
    \vspace{0.4em}
    {\large\heiti\bfseries 面向数据与标签污染的统一评测大底座与主动去噪策略研究 \par}
    \vspace{1.5em}
    {\large 杨景翔\ (2351576) \quad 徐云鹏\ (2353583) \quad 陈艺龙\ (2352359) \quad 李雪菲\ (2354093) \par}
    \vspace{0.5em}
    {\small 数据分析与数据挖掘期末项目 \par}
    \vspace{2em}
\end{center}

\begin{quote}
    \small
    {\heiti\bfseries 摘要：} 异常检测（Anomaly Detection）作为数据挖掘与智能诊断系统的基石，旨在海量、高维的数据流中捕获稀有、偏离常规的少量异常样本。然而，在实际工程落地中，传统的异常检测评估面临着“跨模态碎片化”和“数据/标签污染”的双重严峻学术与工业痛点。为了攻克这一问题，本文构建并实现了一个\textbf{大一统的跨模态异常检测抗污染基准平台}，跨越表格、时序、图（Graph）、图像（CV）和文本（NLP）五大核心模态。
    
    我们系统评估了 14 种经典与前沿算法在 29 个真实异构数据集上的表现，并进行了四组大规模阶梯式实验。首先，我们横向对比了算法在完全干净状态下的性能上限（Exp1）；其次，描绘了在 0\% 至 20\% 的噪声梯度下，算法性能衰退的渐进性退化曲线（Exp2）；然后，通过归一化难度分析，探究了异构模态的分类面搜索硬度（Exp3）；最后，针对重度对称标签翻转噪声，设计并部署了一种基于无监督去噪的即插即用型主动鲁棒防御包装器（Exp4），消融评估了样本剪裁（Trim）和标签翻转（Flip）策略。
    
    大规模评测结果表明：我们提出的以无监督重构（如自编码器、隔离森林）和空间分布（如 IQR）为清洁器核心的 \texttt{IQR\_Trim} 防御策略表现出最强大的抗噪保护作用，在 20\% 极端污染下近乎无损干净数据表现，同时使多层感知机（MLP）和强树集成模型的 AUC-ROC 绝对提升高达 \textbf{2.52\%}；相比之下，标签翻转策略容易引入高维二次噪声，其性能在各维消融中全线弱于剪裁。最后，我们剖析并界定了该防御 wrapper 的失效边界：由于剪裁防御容易破坏 Transformer 上下文和弱线性边界支持点，导致其在元学习大模型 TabPFN 以及线性逻辑回归（LR）中起到了负作用（即满足 \emph{No Free Lunch} 定律）。本研究不仅为跨模态算法鲁棒评估打下了统一的学术底座，更对高鲁棒工业级异常检测系统的实际部署提供了极具含金量的理论指导。
    
    \vspace{0.5em}
    {\heiti\bfseries 关键词：} 异常检测； 跨模态基准； 标签污染； 鲁棒学习； 主动防御； 失效边界
\end{quote}

\vspace{1.5em}

